%% =====================================================================
%%  T-08-02  The Engagement Decision
%%  -------------------------------------------------------------------
%%  One idea per file. Core is mandatory. Slots in canonical order.
%%  Max 700 words. No layout commands. No filler. New source material
%%  for this entry goes in sources/<BOOK>/T-08-02.tex -- never by
%%  rewriting this file (docs/RULES.md R0.1).
%% =====================================================================
%% @kind: topic
%% @id: T-08-02
%% @title: The Engagement Decision
%% @subtitle: Whether to play at all, made before the play
%% @chapter: c08
%% @order: 20
%% @status: stable
%% @sources: B1, B3
%% @updated: 2026-09-19
%% @allow-rewrite: slot migration to the seven-question format (docs/RULES.md section 2)

\topic{T-08-02}{The Engagement Decision}{Whether to play at all, made before the play}


\begin{core}
Before technique, before information, before framing, there is one decision: whether to engage this person about this thing at this time. It is the decision that determines most outcomes, and it is the one most often skipped because it happens before the conversation rather than inside it.
\end{core}

\begin{mechanism}
Engagement is not neutral. Entering an exchange commits attention, reveals interest, opens a relationship that will need maintaining, and --- if the counterpart is of the type described in T-08-01 --- potentially opens an account that will be settled later. All of those costs are incurred whether or not you win. People skip the decision because it is not framed as one: the exchange arrives and they are in it. Making it explicit converts a default into a choice, and a choice can be declined at near-zero cost, which no later move can.
\end{mechanism}

\begin{application}
  \item State the objective before the exchange. No objective, no engagement.
  \item Estimate the upside, the downside, and the cost of the relationship afterwards.
  \item Type the counterpart using T-03-01 and T-08-01 before deciding.
  \item Check whether the timing is yours or theirs. An exchange on their schedule is already partly lost.
  \item Where you decline, decline without explanation and without insult --- both create the account you were avoiding.
  \item Where you engage, engage fully. Half-engagement pays the costs and forfeits the benefits.
\end{application}

\begin{feedback}
  \item You are engaged because the other person initiated, not because you decided.
  \item You cannot state what you want from the exchange.
  \item The upside is small and the downside is open-ended.
  \item You are engaged out of curiosity, irritation, or an inability to let something stand.
  \item Winning would change nothing material.
\end{feedback}

\begin{failure}
Declining has costs that accumulate quietly: reputations for passivity, missed positions, and a pattern of conceding ground that others learn to exploit. Both sources note that non-engagement is itself a move with a payoff, and that payoff is not always positive. There is also a practical floor --- some engagements are unavoidable, involving a manager, a relative, an institution --- and for those the decision reduces to choosing the ground and limiting the exposure rather than declining.
\end{failure}

\begin{countermeasures}
  \item Reserve the right not to play. It is available far more often than it feels like, and the feeling that you must respond is usually manufactured.
  \item Answer a provocation with nothing. Attention withheld costs the provoker more than any reply, and costs you no account.
  \item Where you must engage, choose the ground and the time rather than accepting theirs.
  \item Keep one alternative always available --- another supplier, another job, another route --- because the ability to walk away is what makes declining real.
  \item Ask what happens if you simply do nothing. The answer is usually better than expected.
\end{countermeasures}

\sources{B1, B3}
\seesources
\seealso{T-03-01, T-08-01, T-21-02, T-24-02}
